\documentclass[10pt,a4paper]{article}

\usepackage[utf8]{inputenc}
\usepackage[T1]{fontenc}
\usepackage[french]{babel}
\usepackage{geometry}
\usepackage{booktabs}
\usepackage{array}
\usepackage{amsmath}
\usepackage{xcolor}
\usepackage{hyperref}
\usepackage{caption}
\usepackage{float}
\usepackage{enumitem}
\usepackage{titlesec}

\geometry{margin=2cm}
\setlength{\parskip}{4pt}
\setlength{\parindent}{0pt}
\setlist{itemsep=2pt, topsep=3pt}
\titlespacing*{\section}{0pt}{10pt}{6pt}
\titlespacing*{\subsection}{0pt}{8pt}{4pt}
\hypersetup{
    colorlinks=true,
    linkcolor=blue,
    urlcolor=blue
}

\title{\textbf{Projet P2M : Détection et visualisation d'intrusions par signaux DAS à l'aide de modèles CNN et LSTM}}
\author{Balsem Zouabi \\ Louay Ayadi}
\date{\today}

\begin{document}

\maketitle

\begin{abstract}
Ce rapport présente le projet P2M, une solution complète de détection, de prévision et de visualisation d'événements
acoustiques à partir de signaux DAS (\textit{Distributed Acoustic Sensing}). L'objectif est d'identifier
automatiquement différents types d'activités autour d'une fibre optique, comme la marche, la course, le
passage d'une voiture ou une manipulation de clôture. La solution proposée combine un modèle CNN 2D pour
la classification des événements, un modèle LSTM pour la prévision temporelle du signal, et une application
web de visualisation destinée au suivi des segments de la fibre. Le CNN atteint une précision globale de
94,45 \% sur l'ensemble de test.
\end{abstract}

\tableofcontents
\newpage

\section{Présentation de la problématique}

Les systèmes DAS permettent de transformer une fibre optique en capteur distribué capable de mesurer des
vibrations le long de son parcours. Cette technologie est intéressante pour la surveillance périmétrique,
la protection d'infrastructures et la détection d'intrusions. Chaque événement produit une signature
spatio-temporelle particulière dans le signal mesuré par la fibre.

La problématique traitée dans ce projet consiste à classifier automatiquement un événement DAS parmi
plusieurs classes, à anticiper l'évolution du signal sur un court horizon temporel et à présenter ces
informations dans une interface lisible pour l'utilisateur. Le système doit distinguer des activités normales
ou environnementales de situations potentiellement liées à une intrusion. Les classes considérées par le
module CNN sont les suivantes :

\begin{center}
\begin{tabular}{lll}
\toprule
\textbf{Classe} & \textbf{Description} & \textbf{Type d'événement} \\
\midrule
car & Passage de voiture & Mobilité \\
construction & Activité de construction & Bruit environnemental \\
fence & Interaction avec une clôture & Intrusion potentielle \\
longboard & Passage en longboard & Mobilité \\
manipulation & Manipulation proche de la fibre & Intrusion potentielle \\
openclose & Ouverture/fermeture & Action mécanique \\
regular & Signal régulier ou normal & État non critique \\
running & Course & Mobilité humaine \\
walk & Marche & Mobilité humaine \\
\bottomrule
\end{tabular}
\end{center}

Le défi principal vient de la nature des signaux DAS : ils sont volumineux, bruités et fortement dépendants
du temps et de la position spatiale le long de la fibre. Une méthode efficace doit donc capturer à la fois
l'information fréquentielle du signal, sa localisation spatiale et son évolution temporelle.

\section{Description de la solution proposée}

\subsection{Architecture globale du projet}

La solution finale est organisée en trois modules complémentaires :

\begin{center}
\begin{tabular}{lll}
\toprule
\textbf{Module} & \textbf{Rôle} & \textbf{Fichiers principaux} \\
\midrule
CNN 2D & Classification des événements DAS & \texttt{CNN-Model/} \\
LSTM & Prévision temporelle du signal futur & \texttt{LSTM-Model/} \\
Application web & Visualisation des alertes et des courbes & \texttt{DAS\_Interface/das-monitor/} \\
\bottomrule
\end{tabular}
\end{center}

Le CNN fournit la classe probable de l'événement détecté. Le LSTM complète cette information par une
prévision de l'évolution du signal, utile pour visualiser la tendance future. L'application web regroupe
ces sorties dans une interface de supervision composée d'une carte de la fibre et de courbes de suivi.

\subsection{Prétraitement des données pour le CNN}

Les données brutes sont stockées sous forme de fichiers HDF5. Pour construire l'ensemble d'apprentissage,
le signal est découpé en fenêtres temporelles. Chaque fenêtre est ensuite transformée en représentation
fréquentielle afin de mieux caractériser les vibrations associées aux événements.

Le pipeline d'extraction utilise les paramètres suivants :

\begin{center}
\begin{tabular}{ll}
\toprule
\textbf{Paramètre} & \textbf{Valeur} \\
\midrule
Taille de fenêtre temporelle & 8192 échantillons \\
Décalage entre fenêtres & 2048 échantillons \\
Sur-échantillonnage FFT & 3 \\
Nombre de coefficients FFT conservés & 1024 \\
Nombre de canaux spatiaux par patch & 17 \\
Forme finale d'un échantillon & $(17, 1024, 1)$ \\
\bottomrule
\end{tabular}
\end{center}

Pour chaque canal central sélectionné, un patch 2D est construit avec 17 canaux voisins. Chaque ligne du
patch correspond à un canal spatial, et chaque colonne correspond à une composante fréquentielle. La
transformation appliquée à chaque signal 1D est :

\begin{enumerate}
    \item centrage du signal par soustraction de la moyenne ;
    \item application d'une fenêtre de Hann ;
    \item calcul de la FFT réelle ;
    \item extraction de la magnitude ;
    \item transformation logarithmique $\log(1+x)$ ;
    \item normalisation par z-score.
\end{enumerate}

Cette représentation permet au CNN de traiter le signal comme une image temps-fréquence/spatio-fréquence.

\subsection{Architecture du modèle CNN}

Le modèle utilisé est un réseau de neurones convolutif 2D. Il reçoit en entrée un patch de forme
$(17, 1024, 1)$ et produit une probabilité pour chacune des 9 classes. L'architecture est composée de
quatre blocs convolutifs suivis de couches pleinement connectées.

\begin{center}
\begin{tabular}{lll}
\toprule
\textbf{Bloc} & \textbf{Couches principales} & \textbf{Sortie} \\
\midrule
1 & Conv2D 32 filtres $(3,7)$ + BatchNorm + MaxPooling & $(8,512,32)$ \\
2 & Conv2D 64 filtres $(3,5)$ + BatchNorm + MaxPooling & $(4,256,64)$ \\
3 & Conv2D 128 filtres $(3,3)$ + BatchNorm + MaxPooling & $(2,128,128)$ \\
4 & Conv2D 256 filtres $(3,3)$ + BatchNorm + MaxPooling & $(1,64,256)$ \\
Classification & Flatten + Dense 256 + Dropout + Dense 9 & 9 classes \\
\bottomrule
\end{tabular}
\end{center}

Le modèle contient 13 795 997 paramètres au total. Parmi eux, 4 598 345 sont entraînables. La fonction de
perte utilisée est l'entropie croisée catégorielle sparse :

\[
\mathcal{L} = - \sum_{i=1}^{N} \log(p_{i,y_i})
\]

où $p_{i,y_i}$ représente la probabilité prédite pour la vraie classe de l'échantillon $i$.
L'optimiseur utilisé est Adam avec un taux d'apprentissage de $10^{-4}$.

\subsection{Organisation de l'entraînement}

L'ensemble extrait contient 13 506 échantillons. Il est divisé en trois parties :

\begin{center}
\begin{tabular}{lll}
\toprule
\textbf{Sous-ensemble} & \textbf{Nombre d'échantillons} & \textbf{Rôle} \\
\midrule
Train & 10 804 & Apprentissage des paramètres \\
Validation & 1 351 & Suivi pendant l'entraînement \\
Test & 1 351 & Évaluation finale \\
\bottomrule
\end{tabular}
\end{center}

Des poids de classes sont utilisés pendant l'entraînement afin de réduire l'effet du déséquilibre potentiel
entre les classes. Le meilleur modèle est sauvegardé sous forme de fichier Keras. Une version plus légère
du modèle, sans l'état de l'optimiseur, est aussi fournie pour l'intégration backend.

\subsection{Module LSTM pour la prévision temporelle}

En complément du CNN, un modèle LSTM a été développé pour prédire l'évolution future du signal DAS. Son rôle
n'est pas de classifier l'événement, mais d'estimer la dynamique temporelle du signal sur un horizon court.
Ce module est utile pour l'interface de visualisation, car il permet d'afficher simultanément le signal actuel
et une prédiction de tendance.

Le prétraitement LSTM commence par l'extraction des canaux actifs à partir des fichiers de description
\texttt{.txt}. Le script de préparation copie ensuite les canaux utiles depuis le fichier HDF5 original vers
un nouveau fichier filtré. Cette étape réduit la taille des données et concentre l'apprentissage sur les zones
actives de la fibre.

\begin{center}
\begin{tabular}{ll}
\toprule
\textbf{Paramètre LSTM} & \textbf{Valeur} \\
\midrule
Longueur de séquence d'entrée & 8192 points \\
Horizon de prédiction & 1024 points \\
Nombre de canaux conservés & 8 \\
Stride & 2048 points \\
Batch size & 32 \\
Nombre maximal d'époques & 15 \\
Modèle sauvegardé & \texttt{best\_lstm\_forecast.keras} \\
\bottomrule
\end{tabular}
\end{center}

L'entrée du LSTM est donc de forme $(8192, 8)$ et la sortie attendue est de forme $(1024, 8)$. La cible
utilisée n'est pas directement le signal brut, mais une enveloppe lissée du signal futur. Cette enveloppe
est obtenue à partir de la valeur absolue du signal suivie d'une moyenne mobile. Ce choix rend la prédiction
plus stable, car l'intensité acoustique est généralement plus prévisible que les oscillations instantanées
du signal brut.

L'architecture combine des couches \texttt{Conv1D}, des normalisations, du \texttt{Dropout}, puis une couche
\texttt{Bidirectional LSTM}. Les convolutions réduisent progressivement la taille temporelle et extraient des
motifs locaux, tandis que le LSTM bidirectionnel capture la dépendance temporelle globale dans la séquence.
La sortie finale est produite par une convolution linéaire qui prédit les 8 canaux futurs.

\subsection{Application web de visualisation}

Une application web a également été ajoutée pour présenter les résultats de manière exploitable. Elle est
développée avec React, \texttt{Recharts} pour les courbes et \texttt{lucide-react} pour les icônes. Elle se
trouve dans le dossier \texttt{DAS\_Interface/das-monitor/}.

L'interface est organisée en trois vues principales :

\begin{itemize}
    \item une vue d'authentification ;
    \item une carte de la fibre DAS avec 166 points correspondant à des segments d'environ 10 mètres ;
    \item une vue détaillée d'un segment, affichant l'état de criticité, l'événement CNN et les courbes LSTM.
\end{itemize}

Dans la carte, les segments sont représentés par des points colorés : vert pour un état normal, orange pour
une alerte et rouge pour un danger. Lorsqu'un segment est sélectionné, l'application affiche l'événement
probable classifié par le CNN et dix diagrammes correspondant à dix sous-sections d'un mètre. Chaque diagramme
compare le signal actuel avec la prédiction temporelle fournie par le LSTM.

Cette interface transforme les sorties des modèles en informations lisibles pour un opérateur : localisation
sur la fibre, niveau de danger et évolution du signal.

\section{Analyse des résultats obtenus}

\subsection{Résultats globaux}

Les performances finales sur l'ensemble de test sont les suivantes :

\begin{center}
\begin{tabular}{ll}
\toprule
\textbf{Métrique} & \textbf{Valeur} \\
\midrule
Loss test & 0,1484 \\
Accuracy test & 0,9445 \\
Accuracy test en pourcentage & 94,45 \% \\
Macro F1-score & 0,9445 \\
Weighted F1-score & 0,9445 \\
\bottomrule
\end{tabular}
\end{center}

Ces résultats montrent que le modèle apprend correctement les signatures des différentes classes. Une
précision de 94,45 \% est élevée pour une tâche multi-classes à 9 catégories, surtout avec des signaux
physiques bruités.

\subsection{Résultats par classe}

\begin{center}
\begin{tabular}{lcccc}
\toprule
\textbf{Classe} & \textbf{Précision} & \textbf{Rappel} & \textbf{F1-score} & \textbf{Support} \\
\midrule
car & 0,9800 & 0,9735 & 0,9767 & 151 \\
construction & 0,9671 & 0,9800 & 0,9735 & 150 \\
fence & 0,9478 & 0,8467 & 0,8944 & 150 \\
longboard & 0,9730 & 0,9600 & 0,9664 & 150 \\
manipulation & 0,9732 & 0,9667 & 0,9699 & 150 \\
openclose & 0,9583 & 0,9200 & 0,9388 & 150 \\
regular & 0,8631 & 0,9667 & 0,9119 & 150 \\
running & 0,8944 & 0,9600 & 0,9260 & 150 \\
walk & 0,9586 & 0,9267 & 0,9424 & 150 \\
\bottomrule
\end{tabular}
\end{center}

Les classes \textit{car}, \textit{construction}, \textit{manipulation} et \textit{longboard} présentent les
meilleurs scores F1, tous supérieurs à 0,96. Cela indique que leurs signatures fréquentielles et spatiales
sont bien séparées par le modèle.

La classe \textit{fence} présente le F1-score le plus faible, avec 0,8944. Son rappel est de 0,8467, ce qui
signifie qu'une partie des événements \textit{fence} est confondue avec d'autres classes. D'après la matrice
de confusion, l'erreur principale concerne la confusion entre \textit{fence} et \textit{regular}. Cette
confusion peut s'expliquer par des signatures de vibration proches lorsque l'interaction avec la clôture est
faible ou lorsque le bruit de fond est important.

\subsection{Matrice de confusion}

La matrice de confusion obtenue sur l'ensemble de test est la suivante. Les lignes représentent les vraies
classes et les colonnes les classes prédites.

\[
\begin{bmatrix}
147 & 1 & 0 & 0 & 0 & 1 & 0 & 2 & 0 \\
0 & 147 & 0 & 1 & 0 & 0 & 0 & 2 & 0 \\
0 & 0 & 127 & 0 & 1 & 0 & 20 & 1 & 1 \\
0 & 1 & 0 & 144 & 0 & 1 & 0 & 4 & 0 \\
1 & 0 & 0 & 0 & 145 & 4 & 0 & 0 & 0 \\
0 & 0 & 2 & 1 & 1 & 138 & 2 & 4 & 2 \\
0 & 0 & 5 & 0 & 0 & 0 & 145 & 0 & 0 \\
1 & 0 & 0 & 2 & 0 & 0 & 0 & 144 & 3 \\
1 & 3 & 0 & 0 & 2 & 0 & 1 & 4 & 139
\end{bmatrix}
\]

L'analyse de cette matrice confirme que la majorité des prédictions se situe sur la diagonale, ce qui
correspond aux classifications correctes. L'erreur la plus marquée est la prédiction de 20 échantillons
\textit{fence} comme \textit{regular}. Les autres confusions restent faibles et dispersées.

\subsection{Discussion}

Les résultats obtenus montrent que l'approche CNN 2D est adaptée à la classification des signaux DAS. Le
choix d'une représentation fréquentielle permet de transformer les signaux bruts en motifs plus facilement
exploitables par des convolutions. L'utilisation de plusieurs canaux spatiaux autour d'un canal central
permet également de tenir compte de la propagation locale des vibrations le long de la fibre.

Cependant, certaines limites restent présentes. La confusion entre \textit{fence} et \textit{regular}
montre que le modèle peut avoir des difficultés lorsque l'événement est proche du bruit de fond ou lorsque
son intensité est faible. De plus, l'évaluation dépend de la qualité des annotations et de la représentativité
des fichiers utilisés pour construire le dataset.

\subsection{Analyse du module LSTM et de l'application}

Le module LSTM apporte une information différente de celle du CNN. Alors que le CNN répond à la question
``quel type d'événement est observé ?'', le LSTM répond à la question ``comment le signal risque-t-il
d'évoluer dans les prochains instants ?''. Cette complémentarité est importante pour une application de
supervision : la classification donne le type d'alerte, tandis que la prévision aide à suivre la dynamique du
signal.

Le modèle LSTM sauvegardé dans \texttt{LSTM-Model/Trained\_Model/best\_lstm\_forecast.keras} montre que le
module de prévision a été entraîné et exporté. Le code prévoit une évaluation par MAE et MSE globales, mais
les valeurs numériques finales de ces métriques ne sont pas présentes dans les fichiers de résultats fournis.
Dans le rapport, l'analyse quantitative détaillée est donc centrée sur le CNN, pour lequel le rapport de
classification et la matrice de confusion sont disponibles.

L'application web constitue la partie de valorisation du projet. Elle permet de visualiser les alertes sur
une topologie de fibre composée de 166 segments, puis d'ouvrir une vue détaillée par segment. Cette vue
combine trois informations : le niveau de criticité, l'événement probable issu du CNN et les courbes de
signal réel/prédit associées au LSTM. Le résultat attendu est une interface exploitable par un utilisateur
non spécialiste du modèle, qui peut localiser rapidement une zone critique et consulter l'évolution du
signal.

\section{Conclusion}

Ce projet a permis de développer une solution complète de traitement de signaux DAS. Le pipeline proposé
combine une extraction de caractéristiques par FFT, un CNN 2D capable de reconnaître 9 classes d'événements,
un LSTM pour la prévision temporelle du signal et une application web de visualisation. Les résultats CNN
montrent une accuracy de 94,45 \% sur l'ensemble de test, ce qui confirme la pertinence de l'approche de
classification.

Pour améliorer le système, plusieurs pistes peuvent être envisagées : enrichir le dataset avec davantage
d'exemples de classes difficiles, appliquer de l'augmentation de données contrôlée, documenter les métriques
finales du LSTM, connecter l'application web à un backend réel, et tester des architectures plus compactes
pour le déploiement.

\section*{Fichiers livrés}

\begin{itemize}
    \item \texttt{CNN-Model/extract\_dataset\_2d.py} : extraction des caractéristiques.
    \item \texttt{CNN-Model/train\_from\_extracted\_2d.py} : entraînement du CNN.
    \item \texttt{CNN-Model/classes\_das\_2d.npy} : noms des classes.
    \item \texttt{best\_cnn\_das\_2d.keras} : modèle complet.
    \item \texttt{best\_cnn\_das\_2d\_backend.keras} : modèle léger pour l'inférence backend.
    \item \texttt{LSTM-Model/Dta\_Preparation/Data\_Preparation.py} : préparation des canaux actifs.
    \item \texttt{LSTM-Model/Model/LSTM.py} : entraînement du modèle de prévision.
    \item \texttt{LSTM-Model/Trained\_Model/best\_lstm\_forecast.keras} : modèle LSTM sauvegardé.
    \item \texttt{DAS\_Interface/das-monitor/} : application web React de visualisation.
\end{itemize}

\end{document}
